\subsection{部分グラフ}

\begin{definition}
  \label{def:CartanMatrix.LowerLabel}
  \lean{CartanMatrix.LowerLabel}
  \leanok
  次のような行列$C_L$を定める：
  \begin{equation}
    (C_L)_{ij} = \begin{cases}
      2, & i = j, \\
      -1, & C_{ij} = -1, \\
      -1, & C_{ij} = -2, \\
      -2, & C_{ij} = -3.
    \end{cases}
  \end{equation}
  $*_L$と$*_k$の操作を組み合わせてできる行列を\textbf{部分グラフ}という．
  Coxeter-Dynkin図形の言葉で言えば，辺の本数を減らす，頂点を減らす操作の組み合わせである．
\end{definition}

正定値の部分グラフもまた正定値である．

\begin{theorem}
  \label{thm:CartanMatrix.sub_of_pos_def}
  \lean{CartanMatrix.sub_of_pos_def}
  \leanok
  \uses{lem:SymmMatrix_leadingSubmatrix_comm, lem:PosDef.submatrix_injective}
  $S(C)$が正定値ならば，$S(C_k)$も正定値である．
\end{theorem}

\begin{proof}
  補題~\ref{lem:SymmMatrix_leadingSubmatrix_comm}より
  \begin{equation}
    S(C_k)
    = (S(C))_k
  \end{equation}
  であり，補題~\ref{lem:PosDef.submatrix_injective}より正定値行列の部分行列は正定値であるから従う．
\end{proof}

\begin{theorem}
  \label{thm:CartanMatrix.sub_of_pos_def'}
  \lean{CartanMatrix.sub_of_pos_def'}
  \leanok
  \uses{}
  $S(C)$が正定値ならば，$S(C_L)$も正定値である．
\end{theorem}

\begin{proof}
  対偶を示す．
  $S(C_L)$が正定値でないとき，$y \in \R^n$で
  \begin{equation}
    y \ne 0, \qquad \sum_{i, j} y_i \cdot (S(C))_{ij} \cdot y_j \le 0
  \end{equation}
  なるものがあることを示せばよい．
  $S(C_L)$が正定値でないから$0 \ne x \in \R^n$で
  \begin{equation}
    \sum_{i, j} x_i \cdot (S(C_L))_{ij} \cdot x_j \le 0
  \end{equation}
  なるものがある．
  \begin{equation}
    y = \transpose{(|x_1|, \ldots, |x_n|)}
  \end{equation}
  とする．

  $y \ne 0$を示す．
  $y = 0$とすると，任意の$i$に対し$y_i = |x_i| = 0$，すなわち$x_i = 0$であるが，これは矛盾する．

  また，
  \begin{align}
    \sum_{i, j} y_i \cdot (S(C))_{ij} \cdot y_j
    &= \sum_{i, j} |x_i| \cdot (S(C))_{ij} \cdot |x_j| \\
    &\le \sum_{i, j} |x_i| \cdot (S(C))_{ij} \cdot |x_j| \\
    &= \sum_{i, j} (S(C_L))_{ij} \cdot |x_i x_j| \\
    &= \sum_{i, j} (S(C_L))_{ij} \cdot x_i x_j \\
    &= \sum_{i, j} x_i \cdot (S(C_L))_{ij} \cdot x_j \\
    &\le 0.
  \end{align}
\end{proof}